%% ma: L10-B13-0108
%% lop: 10
%% chuong: 5
%% bai: 13
%% dang: 10.13.13
%% loai: DS
%% muc: [NB, NB, TH, TH]
%% dap_an: "SSĐĐ"
%% nguon: "chuyen-de-cac-so-dac-trung-cua-mau-so-lieu-khong-ghep-nhom-toan-10-chuong-trinh-moi.pdf (GV. Phan Nhật Linh) (P03), Bài 02 Dạng 2 Phần II Câu 12"
%% kien_thuc: "Lớp 10 Bài 13 (số trung bình, trung vị, tứ phân vị, mốt)"
%% trang_thai: chinh-thuc
%% ghi_chu: "Tự duyệt theo yêu cầu giáo viên 22/10/2026."
\begin{ex}
Kết quả điểm kiểm tra môn Toán của $40$ học sinh lớp 10A được trình bày ở bảng sau:
\begin{center}
\begin{tabular}{|c|c|c|c|c|c|c|c|c|}\hline
Điểm & $4$ & $5$ & $6$ & $7$ & $8$ & $9$ & $10$ & Cộng\\\hline
Tần số & $2$ & $8$ & $7$ & $10$ & $8$ & $3$ & $2$ & $40$\\\hline
\end{tabular}
\end{center}
Xét tính đúng sai của các khẳng định sau:
\choiceTF{Mốt của mẫu số liệu là $M_o=10$}{Số học sinh có điểm từ $8$ trở lên là $5$}{\True Số trung bình của mẫu số liệu (làm tròn đến hàng phần mười) là $6{,}8$}{\True Trung vị của mẫu số liệu là $7$}
\loigiai{a) Sai: $M_o=7$ (tần số $10$).

b) Sai: $8+3+2=13$ học sinh.

c) Đúng: $\overline x=\dfrac{271}{40}=6{,}775\approx6{,}8$.

d) Đúng: tần số tích lũy $2;10;17;27$ nên $x_{20}=x_{21}=7$, $M_e=7$.}
\end{ex}
